% =============================================================================
% Chapter: Methods --- Forecasting Models and Evaluation (modelling thread)
% Source of pinned setup values: configs/*.py on this branch +
% docs/superpowers/specs/2026-06-04-final-model-experiments-and-report-design.md
% Numeric DNN hyperparameters are the committed (interim) tuned values and are
% refreshed by the paper-exogenous re-tune that accompanies the final benchmark.
% =============================================================================

\chapter{Methods --- Forecasting Models and Evaluation}
\label{ch:methods-models}

This chapter specifies the forecasting models that the platform serves and the
protocol by which they are compared. Three model families are studied: the
Lasso Estimated AutoRegressive ensemble (LEAR) and the deep neural network
(DNN) of Lago et al.~\cite{lago2021}, and the encoder-only Transformer of
Llorente and Portela~\cite{llorente2025transformer}. Each is reimplemented
faithfully and run on recent day-ahead data for two European bidding zones,
DK1 (Nord Pool) and DE (EPEX). The chapter first states what "faithful
replication" means and why the comparison is framed in relative rather than
absolute terms (\S\ref{sec:mm-framing}), then describes the data and feature
construction (\S\ref{sec:mm-data}), the three models as implemented
(\S\ref{sec:mm-models}), the rolling recalibration and backtesting scheme
(\S\ref{sec:mm-recalibration}), the evaluation metrics and significance tests
(\S\ref{sec:mm-eval}), the hyperparameter-tuning protocol
(\S\ref{sec:mm-tuning}), and finally the documented deviations from the source
papers (\S\ref{sec:mm-deviations}).

% -----------------------------------------------------------------------------
\section{Replication framing}
\label{sec:mm-framing}

The error magnitudes reported in the source papers were obtained on 2016--2018
test periods whose price level and volatility differ sharply from the
2025--2026 regime studied here; an absolute mean-error figure is therefore not
comparable across the two eras. The comparison this work makes is consequently
threefold. First, \emph{methodology replication}: the same model
specifications, feature sets, recalibration scheme, metrics, and statistical
tests as the papers, executed on current data. Second, \emph{relative
accuracy}: the relative mean absolute error (rMAE), which Lago et
al.~\cite{lago2021} designed precisely for cross-period and cross-market
comparison, together with the ordinal ranking of the models against the
rankings the papers report. Third, \emph{head-to-head significance}: paired
Diebold--Mariano~\cite{diebold1995dm} and Giacomini--White~\cite{giacomini2006gw}
tests across the three models and a weekly-naive baseline, per market, on
identical test days. Because all six experimental runs (three models $\times$
two markets) share a single test window, the paired tests are valid: every pair
of models is compared on the same delivery days.

% -----------------------------------------------------------------------------
\section{Data and features}
\label{sec:mm-data}

\begin{table}[h]
  \centering
  \caption{Datasets. Both markets span 2022-03-10 to 2026-05-31 (ENTSO-E Transparency Platform).}
  \label{tab:mm-datasets}
  \begin{tabularx}{\linewidth}{@{}l l l X@{}}
    \toprule
    \textbf{Market} & \textbf{Resolution} & \textbf{Exogenous (paper-exact)} & \textbf{Paper analogue} \\
    \midrule
    DK1 & hourly (24)\textsuperscript{a} & TotalLoad + Wind (on+offshore)        & Nord Pool: load + wind~\cite{lago2021} \\
    DE  & hourly (24)\textsuperscript{a} & TotalLoad + Renewables (wind+solar)   & EPEX-DE: load + wind\&solar~\cite{lago2021} \\
    \\[2pt] \multicolumn{4}{@{}p{\linewidth}@{}}{\footnotesize\textsuperscript{a} Both markets clear in 15-min blocks (DE throughout; DK1 since 2025-10-01); prices and exogenous forecasts are downsampled to an arithmetic hourly mean (\texttt{resample\_to\_hourly}).} \\
    \bottomrule
  \end{tabularx}
\end{table}

All data are obtained from the ENTSO-E Transparency Platform and stored in
ClickHouse, from which the preprocessor loads a price series and a configurable
set of exogenous forecast series for each market. Following Lago et
al.~\cite{lago2021}, Eq.~3, the feature vector for each delivery day combines
autoregressive price lags at $d{-}1$, $d{-}2$, $d{-}3$, and $d{-}7$; exogenous
forecasts at $d0$, $d{-}1$, and $d{-}7$; and seven day-of-week dummies. For a
day discretised into $n$ delivery periods this yields
\begin{equation}
  \underbrace{4n}_{\text{price lags}} + \underbrace{2 \cdot 3n}_{\text{two exogenous} \times \text{3 lags}} + \underbrace{7}_{\text{weekday}}
  \label{eq:mm-featurecount}
\end{equation}
features, i.e.\ $247$ at hourly resolution ($n{=}24$) and $967$ at quarter-hourly
resolution ($n{=}96$). The exogenous set is restricted, per market, to exactly
the two signals the papers use: for DK1 the zonal load forecast and an
aggregated wind forecast (onshore plus offshore); for DE the zonal load forecast
and an aggregated renewables forecast (onshore plus offshore wind plus solar).
The aggregates are formed at load time by summing the raw ENTSO-E curves: the
member curves are inner-joined on their timestamps before summing so that a
partial sum can never occur, and all exogenous columns are then inner-joined
together, so the coarsest available series anchors the common time grid.

The target undergoes the asinh variance-stabilising transform of Uniejewski et
al.~\cite{uniejewski2018vst}, which compresses the heavy tails and price spikes
characteristic of day-ahead series while, unlike a logarithm, remaining defined
for the negative prices that occur under high renewable infeed. Both markets
now clear in fifteen-minute blocks --- DE throughout the sample and DK1 since
2025-10-01 --- but prices and exogenous forecasts are downsampled to an
arithmetic hourly mean (24 periods per day) before features are constructed.
This keeps the two markets comparable with each other and with the hourly
reference benchmarks of Lago et al.~\cite{lago2021} and Llorente and
Portela~\cite{llorente2025transformer}; modelling DK1 natively at 96 periods is
in any case infeasible, as it has under one month of quarter-hourly history
before the test window. Finally, ENTSO-E lacks DK1 wind forecasts for a small
number of days inside the sample; affected days are
dropped identically from every DK1 model, so the paired significance tests
remain valid.

% -----------------------------------------------------------------------------
\section{Models}
\label{sec:mm-models}

Table~\ref{tab:mm-models} summarises the three model setups, which are identical
across both markets except for the market-specific exogenous aggregate and the
per-market DNN hyperparameters.

\begin{table}[h]
  \centering
  \caption{Model setups (one row per model; identical across both markets unless noted).}
  \label{tab:mm-models}
  \begin{tabularx}{\linewidth}{@{}l X X l@{}}
    \toprule
    \textbf{Model} & \textbf{Specification} & \textbf{Calibration} & \textbf{Refit} \\
    \midrule
    LEAR        & One Lasso per delivery period; $\lambda$ by \texttt{LassoCV} (time-series CV) & ensemble over 56 / 84 / 1092\,d & daily \\
    DNN         & MLP; architecture HPs by Optuna TPE (ex-ante validation); 4-seed ensemble at serving, single seed in backtest & 730\,d & daily \\
    Transformer & Encoder stack per~\cite{llorente2025transformer}; paper Nord Pool hyperparameters & 1092\,d & weekly \\
    \bottomrule
  \end{tabularx}
\end{table}

\subsection{LEAR}

LEAR is a direct multi-period regressor: one independent Lasso model is fitted
per delivery period, so an hourly market trains 24 separate Lasso estimators,
each predicting a single hour from the full feature vector of
Eq.~\eqref{eq:mm-featurecount}. The L1 penalty performs simultaneous shrinkage
and feature selection, which is what makes the parameter-rich input space
tractable. The regularisation strength $\lambda$ is selected automatically per
period by \texttt{LassoCV} with a five-fold \texttt{TimeSeriesSplit}
cross-validation (maximum 5000 coordinate-descent iterations), the configuration
denoted \texttt{alpha="auto"}. LEAR carries no learned ensembling of its own;
the ensemble averaging it benefits from is supplied by the multi-window
recalibration scheme of \S\ref{sec:mm-recalibration}.

\subsection{Deep neural network}

The DNN is a multi-layer perceptron with two hidden layers that maps the feature
vector to all $n$ delivery periods in a single forward pass (a
multiple-input--multiple-output design). It is trained to minimise the mean
absolute error with the Adam optimiser, with optional L1 weight regularisation,
optional batch normalisation, and early stopping (the best weights, by
validation loss, are restored from memory rather than from disk). The serving
prediction is the arithmetic mean of four replicates trained from distinct
random seeds, following the ensemble recommendation of Lago et
al.~\cite{lago2021}, \S4.4. This four-seed averaging is applied only on the
production serving path; the rolling backtest that produces the
Chapter~\ref{ch:results-models} results trains a single network per
recalibration at a fixed seed, so the reported DNN metrics carry the seed
variance of one draw (revisited in \S\ref{sec:rm-validity}). The six architecture hyperparameters --- hidden-layer
width, activation, dropout, learning rate, the batch-normalisation flag, and the
L1 penalty --- are selected per market by the tuning protocol of
\S\ref{sec:mm-tuning}; Table~\ref{tab:mm-dnn-hp} lists the values in force.

\begin{table}[h]
  \centering
  \caption{Tuned DNN hyperparameters per market. Selected by the protocol of \S\ref{sec:mm-tuning} on the ex-ante validation window; the architecture-only search is re-run whenever the feature set changes, so these values are refreshed by the paper-exogenous re-tune that accompanies the final benchmark.}
  \label{tab:mm-dnn-hp}
  \begin{tabularx}{\linewidth}{@{}l l l l l l X@{}}
    \toprule
    \textbf{Market} & \textbf{Width} & \textbf{Activation} & \textbf{Dropout} & \textbf{Learning rate} & \textbf{Batch norm} & \textbf{L1 penalty} \\
    \midrule
    DK1 & 191 & tanh & 0.203 & $4.8\times10^{-4}$ & no  & $6.8\times10^{-8}$ \\
    DE  & 373 & relu & 0.204 & $6.9\times10^{-4}$ & yes & $7.1\times10^{-7}$ \\
    \bottomrule
  \end{tabularx}
\end{table}

\subsection{Transformer}

The Transformer follows Llorente and Portela~\cite{llorente2025transformer} and
is encoder-only with two parallel input branches. The first branch takes the
past-price sequence of the previous $m{=}14$ days, each day a vector of $n$
period prices; each day is linearly embedded to a 128-dimensional vector, a
learnable positional encoding is added, and the sequence is passed through a
stack of six Transformer encoder layers (eight attention heads, feed-forward
dimension 1024, dropout 0.2), of which the final position is taken as the price
representation. The second branch linearly embeds the day-ahead exogenous
forecasts. The two representations are concatenated and mapped by a
single-hidden-layer perceptron (width 128) to the $n$ output periods. The model
is trained with Adam at a learning rate of $10^{-4}$, batch size 32, for up to
200 epochs with early stopping (patience 25) on an internal validation split,
under a step learning-rate schedule. The number of exogenous inputs is fixed at
two to match the paper specification, and the remaining hyperparameters follow
the paper's reported architecture and search ranges rather than a per-market
optimum --- the 128-dimensional embedding lies within the paper's search range
but below its Nord Pool optimum of 512. Per-market Transformer tuning is left as
future work (\S\ref{sec:mm-deviations}). Because the wide feature vector this
model consumes is interpreted positionally rather than by column name, it is
produced by a dedicated feature pipeline
(\texttt{default\_transformers\_for\_transformer}) that lays out the
$m{\times}n$ price block followed by the exogenous block in a fixed order.

% -----------------------------------------------------------------------------
\section{Recalibration and backtesting}
\label{sec:mm-recalibration}

All three models are evaluated by a single walk-forward backtesting engine that
reproduces the rolling-recalibration scheme of Lago et al.~\cite{lago2021}. For
each delivery day in the test window, and for each calibration-window length in
the model's ensemble, the engine slices the training data to the most recent
window of days strictly preceding the target day, fits the feature and target
scalers on that slice alone, trains the model, and predicts the target day; the
day's forecast is the arithmetic mean of the predictions across the
calibration windows. To bound compute, each window keeps a cached fitted model
and is refitted only when the number of days since its last fit reaches the
configured refit interval --- daily for LEAR and the DNN, weekly for the
Transformer, following the respective papers. LEAR uses the paper ensemble of
56-, 84-, and 1092-day windows (the 1456-day window of the original
specification is dropped because it predates the available data); the DNN uses a
single 730-day window; and the Transformer uses a single 1092-day window, which
was verified to leave enough usable sequence rows before the test start in both
markets (at least 1092 in both; 1265 measured for DE). The train/test split is
date-based: training data are all rows strictly before the test start date, and
no validation holdout is carved out of the backtest itself (the validation
window of \S\ref{sec:mm-tuning} is used only by the tuning command).

% -----------------------------------------------------------------------------
\section{Evaluation protocol}
\label{sec:mm-eval}

The test window runs from 2025-10-24 to 2026-05-31 (220 calendar days) and is
identical for all six runs. Four accuracy metrics are reported. For $N$ aligned
prediction--actual pairs $(\hat{y}_i, y_i)$,
\begin{align}
  \mathrm{MAE}   &= \frac{1}{N}\sum_{i=1}^{N} \lvert y_i - \hat{y}_i \rvert,
  & \mathrm{RMSE} &= \sqrt{\frac{1}{N}\sum_{i=1}^{N} (y_i - \hat{y}_i)^2}, \\
  \mathrm{sMAPE} &= \frac{100}{N}\sum_{i=1}^{N} \frac{\lvert y_i - \hat{y}_i\rvert}{\tfrac{1}{2}(\lvert y_i\rvert + \lvert \hat{y}_i\rvert)},
  & \mathrm{rMAE} &= \frac{\mathrm{MAE}(\mathbf{y}, \hat{\mathbf{y}})}{\mathrm{MAE}(\mathbf{y}, \hat{\mathbf{y}}^{\text{naive}})},
\end{align}
where the naive benchmark $\hat{\mathbf{y}}^{\text{naive}}$ is the weekly-naive
forecast (the realised price exactly seven days earlier), as recommended by Lago
et al.~\cite{lago2021}, \S5.4. The rMAE is the headline metric because, being
normalised by the same-market naive error, it is the figure that can be compared
across periods and markets.

Statistical significance is assessed with the multivariate Giacomini--White
test~\cite{giacomini2006gw, lago2021}. For each ordered pair of models $(A,B)$
the daily loss differential
\begin{equation}
  d_t = \sum_{h=1}^{n} \bigl(\lvert e^{A}_{t,h}\rvert - \lvert e^{B}_{t,h}\rvert\bigr)
\end{equation}
is formed over the $n$ delivery periods of day $t$ (the $L^1$ norm of the
errors), and the test statistic $\bar{d}/\hat{\sigma}_{\bar d}$ is computed with
a Newey--West heteroskedasticity- and autocorrelation-consistent standard error
using a lag of $\lfloor T^{1/4}\rfloor$ over the $T$ test days. The univariate
Diebold--Mariano test~\cite{diebold1995dm} is reported as a complement. Models
are compared on the test days that all of them produced, since the different
feature pipelines drop gap and daylight-saving days slightly differently across
model families. Finally, computation cost is reported as the training time per
recalibration, read from the experiment tracker, as the analogue of Lago et
al.~\cite{lago2021}, Table~4.

% -----------------------------------------------------------------------------
\section{Hyperparameter tuning}
\label{sec:mm-tuning}

The DNN is the only model whose hyperparameters are searched; LEAR selects
$\lambda$ internally by cross-validation, and the Transformer reuses the paper
optimum. Tuning uses Optuna's Tree-structured Parzen Estimator over a six-dimensional
architecture-only search space: hidden-layer width (16--512, log scale),
activation (one of relu, tanh, elu, softplus, sigmoid), dropout (0--0.5),
learning rate ($10^{-5}$--$10^{-1}$, log scale), the batch-normalisation flag,
and the L1 penalty ($10^{-8}$--$10^{-2}$, log scale). The objective is the mean
MAE over the validation window of 2025-01-01 to 2025-10-23 --- the 42 weeks
immediately preceding the test window, kept strictly ex-ante per Lago et
al.~\cite{lago2021}, \S4.3.2 --- evaluated under the same walk-forward engine
used for the backtest. Every trial is persisted to the
\texttt{algo.hyperparam\_trials} table so that a search resumes across sessions.
The ten binary feature- and selector-toggles of Lago et al.\ (\S4.3.3), together
with their \texttt{preprocessing} and \texttt{init\_method} choices, are
deliberately omitted: the preprocessor and model constructors do not expose the
corresponding switches, so including them in the search space would create inert
dimensions that consume the trial budget without affecting the model.

% -----------------------------------------------------------------------------
\section{Deviations from the source papers}
\label{sec:mm-deviations}

The replication is faithful in specification but departs from the source papers
in a small number of documented respects, summarised in
Table~\ref{tab:mm-deviations}. Each deviation is driven either by the length of
the available data history (which begins on 2022-03-10) or by a compute budget,
and none changes the model class or the feature definition.

\begin{table}[h]
  \centering
  \caption{Documented deviations from~\cite{lago2021} and~\cite{llorente2025transformer}.}
  \label{tab:mm-deviations}
  \begin{tabularx}{\linewidth}{@{}l X X@{}}
    \toprule
    \textbf{Aspect} & \textbf{Paper} & \textbf{This work (reason)} \\
    \midrule
    Test length            & $\geq$ 1--2 years            & 31 weeks (price history starts 2022-03; validation and calibration consume the rest) \\
    DNN tuning budget      & 1500 TPE iterations          & 30 trials (compute; the optimum is stable early in prior studies) \\
    DNN search space       & architecture + feature selectors & architecture only (feature set fixed to the paper specification) \\
    LEAR windows           & 56 / 84 / 1092 / 1456\,d     & 56 / 84 / 1092\,d (1456 predates the data) \\
    DNN window             & 4 years                      & 730\,d (data and compute constraint) \\
    $\lambda$ selection    & LARS+AIC, then coordinate-descent refit & \texttt{LassoCV} time-series CV (standard; close accuracy) \\
    Market resolution      & hourly (paper era)           & both markets now 15-min; downsampled to an hourly mean for cross-market and cross-paper comparability \\
    Transformer pipeline   & train-once + renormalise     & calibration window + weekly refit (consistency with the other models) \\
    Transformer HPs        & tuned per dataset            & paper architecture / search-range defaults reused (per-market tuning is future work) \\
    \bottomrule
  \end{tabularx}
\end{table}
